Mapa pokrytí kolektivními smlouvami (OECD/ICTWSS) v~EU27 odhaluje dramatický geografický kontrast: zatímco Rakousko a~Francie dosahují \SIrange{95}{98}{\percent} pokrytí díky erga omnes mechanismu nebo právní tradici rozšiřování \ac{KS}, ČR se pohybuje kolem \SI{35}{\percent} --- tedy méně než třetina pracovníků je pokryta \ac{KS}.
Tato propast není výsledkem nižší odborové tradice ani slabší ekonomiky: je výsledkem vědomé institucionální volby --- absence zákonného mechanismu pro rozšíření \ac{KS} na celé odvětví.
Ve státech s~vysokým pokrytím funguje erga omnes (FR, AT) nebo kvazipovinná registrace odvětvových smluv (DE, NL): i~zaměstnanec, který není členem odborů, je automaticky pokryt sjednanoukategorizací a~tarifní tabulkou --- čímž odpadá motiv vyhýbat se odborové příslušnosti.
Zavedení analogického mechanismu v~ČR --- tj.~zákonné rozšíření odvětvových \ac{KS} na všechny zaměstnavatele v~daném odvětví (§~7 zákona o~kolektivním vyjednávání~\cite{zakon_kv_1991} ve zreformované podobě) --- je ústřední institucionální reformou, k~níž tato práce dospívá jako ke klíčovému doporučení.

Mapa pokrytí kolektivními smlouvami (OECD/ICTWSS) v~EU27 odhaluje dramatický geografický kontrast: zatímco Rakousko a~Francie dosahují \SIrange{95}{98}{\percent} pokrytí díky erga omnes mechanismu nebo právní tradici rozšiřování \ac{KS}, ČR se pohybuje kolem \SI{35}{\percent} --- tedy méně než třetina pracovníků je pokryta \ac{KS}.
Tato propast není výsledkem nižší odborové tradice ani slabší ekonomiky: je výsledkem vědomé institucionální volby --- absence zákonného mechanismu pro rozšíření \ac{KS} na celé odvětví.
Ve státech s~vysokým pokrytím funguje erga omnes (FR, AT) nebo kvazipovinná registrace odvětvových smluv (DE, NL): i~zaměstnanec, který není členem odborů, je automaticky pokryt sjednanoukategorizací a~tarifní tabulkou --- čímž odpadá motiv vyhýbat se odborové příslušnosti.
Zavedení analogického mechanismu v~ČR --- tj.~zákonné rozšíření odvětvových \ac{KS} na všechny zaměstnavatele v~daném odvětví (§~7 zákona o~kolektivním vyjednávání~\cite{zakon_kv_1991} ve zreformované podobě) --- je ústřední institucionální reformou, k~níž tato práce dospívá jako ke klíčovému doporučení.

Mapa pokrytí kolektivními smlouvami (OECD/ICTWSS) v~EU27 odhaluje dramatický geografický kontrast: zatímco Rakousko a~Francie dosahují \SIrange{95}{98}{\percent} pokrytí díky erga omnes mechanismu nebo právní tradici rozšiřování \ac{KS}, ČR se pohybuje kolem \SI{35}{\percent} --- tedy méně než třetina pracovníků je pokryta \ac{KS}.
Tato propast není výsledkem nižší odborové tradice ani slabší ekonomiky: je výsledkem vědomé institucionální volby --- absence zákonného mechanismu pro rozšíření \ac{KS} na celé odvětví.
Ve státech s~vysokým pokrytím funguje erga omnes (FR, AT) nebo kvazipovinná registrace odvětvových smluv (DE, NL): i~zaměstnanec, který není členem odborů, je automaticky pokryt sjednanoukategorizací a~tarifní tabulkou --- čímž odpadá motiv vyhýbat se odborové příslušnosti.
Zavedení analogického mechanismu v~ČR --- tj.~zákonné rozšíření odvětvových \ac{KS} na všechny zaměstnavatele v~daném odvětví (§~7 zákona o~kolektivním vyjednávání~\cite{zakon_kv_1991} ve zreformované podobě) --- je ústřední institucionální reformou, k~níž tato práce dospívá jako ke klíčovému doporučení.

Mapa pokrytí kolektivními smlouvami (OECD/ICTWSS) v~EU27 odhaluje dramatický geografický kontrast: zatímco Rakousko a~Francie dosahují \SIrange{95}{98}{\percent} pokrytí díky erga omnes mechanismu nebo právní tradici rozšiřování \ac{KS}, ČR se pohybuje kolem \SI{35}{\percent} --- tedy méně než třetina pracovníků je pokryta \ac{KS}.
Tato propast není výsledkem nižší odborové tradice ani slabší ekonomiky: je výsledkem vědomé institucionální volby --- absence zákonného mechanismu pro rozšíření \ac{KS} na celé odvětví.
Ve státech s~vysokým pokrytím funguje erga omnes (FR, AT) nebo kvazipovinná registrace odvětvových smluv (DE, NL): i~zaměstnanec, který není členem odborů, je automaticky pokryt sjednanoukategorizací a~tarifní tabulkou --- čímž odpadá motiv vyhýbat se odborové příslušnosti.
Zavedení analogického mechanismu v~ČR --- tj.~zákonné rozšíření odvětvových \ac{KS} na všechny zaměstnavatele v~daném odvětví (§~7 zákona o~kolektivním vyjednávání~\cite{zakon_kv_1991} ve zreformované podobě) --- je ústřední institucionální reformou, k~níž tato práce dospívá jako ke klíčovému doporučení.

\input{texparts/python/kv_coverage_map}



